\chapter{Trabalhos Relacionados}
\label{sec:RW}
    O Algoritmo Genético com Fertilização In Vitro (IVF/GA)~\cite{camilo-junior2011} é uma extensão do algoritmo genético canônico que realiza busca local através de um processo denominado Fertilização In Vitro (IVF).

    Um estudo recente de Sampaio~\cite{sampaio2017ivf} apresentou uma abordagem diferente para o método IVF em otimização multiobjetivo. O estudo introduziu uma estratégia distinta para coletar indivíduos parentais iniciais, aprimorando a diversidade da população. No entanto, a estratégia de manipulação genética permaneceu inalterada, com apenas a melhor descendência substituindo o parental atual. Esta abordagem mantém os princípios fundamentais do método original e oferece oportunidades para a qualidade da solução.

    Publicações como IVF/NSGA-II~\cite{sampaio2017ivf} e IVF/GDE3~\cite{sampaio2019ivf} destacaram as contribuições positivas do acoplamento do método IVF com algoritmos evolutivos e de evolução diferencial populares como NSGA-II e GDE3. Estes trabalhos citados acima demonstraram o potencial do módulo para aprimorar as capacidades de convergência e diversificação do algoritmo hospedeiro. O acoplamento do método IVF com tais algoritmos fornece uma estrutura robusta para resolver problemas de otimização complexos e equilibrar o trade-off exploração-explotação.

    Além da hibridização direta com módulos de busca local como IVF, outros aprimoramentos significativos para MOEAs, incluindo SPEA2, têm se concentrado em adaptá-los a classes específicas de problemas desafiadores ou melhorar seus mecanismos centrais. Por exemplo, Problemas de Otimização Multiobjetivo com Restrições (CMOPs) apresentam consideráveis desafios quanto à satisfação de restrições e ao equilíbrio entre diversidade populacional e convergência. Para abordar isso, foi introduzido o framework de Otimização Multiforme (MFO), no qual o CMOP original é tratado juntamente com uma tarefa CMOP auxiliar em um cenário de otimização multitarefa~\cite{jiao2023multiform}. Quando implementado com SPEA2 (MFO-SPEA2), este framework não modifica intrinsecamente o algoritmo SPEA2 central, mas o integra em uma arquitetura de dupla tarefa adequada para CMOPs. A tarefa auxiliar, derivada do CMOP primário, opera com restrições inicialmente relaxadas que são dinamicamente e progressivamente restringidas. A transferência de conhecimento entre essas tarefas, uma forma implícita de hibridização de informações de busca, utiliza mecanismos de acasalamento e seleção personalizados para cada formulação de problema para guiar a busca evolutiva em direção a soluções Pareto-ótimas, aproveitando informações de domínios de soluções viáveis e inviáveis~\cite{jiao2023multiform}. Esta abordagem permitiu ao MFO-SPEA2 mostrar desempenho superior contra métodos representativos de tratamento de restrições e outros CMOEAs estado-da-arte em suítes de teste padrão e um problema real de síntese de arranjo de antenas~\cite{jiao2023multiform}.

    Dentro da otimização multi objetivos, algoritmos baseados em Pareto podem ter dificuldades devido à eficácia reduzida da relação de dominância de Pareto. Isso pode levar mecanismos de manutenção de diversidade a inadvertidamente guiar a busca para longe da verdadeira frente de Pareto~\cite{li2014shift}. Para abordar isso, a estratégia de Estimação de Densidade Baseada em Deslocamento (SDE) foi introduzida como uma melhoria ao componente de preservação de diversidade dos MOEAs~\cite{li2014shift}. SPEA2+SDE implementa esta estratégia modulando a estimação de densidade do k-ésimo vizinho mais próximo do SPEA2, efetivamente hibridizando o cálculo de densidade com informações de convergência~\cite{li2014shift}. SDE incorpora informações tanto distribucionais quanto relacionadas à convergência; durante a avaliação de densidade para um indivíduo focal, SDE computacionalmente ``desloca'' as posições de outros indivíduos baseado em suas características de convergência relativas ao indivíduo focal. Assim, soluções mal convergidas são percebidas como estando em regiões mais densas, diminuindo sua probabilidade de seleção~\cite{li2014shift}. Avaliações empíricas mostraram que SPEA2+SDE supera marcadamente o desempenho do SPEA2 canônico em problemas de otimização de muitos objetivos, especialmente conforme a dimensionalidade do espaço objetivo aumenta. Também demonstrou competitividade robusta contra outros algoritmos líderes projetados para cenários de muitos objetivos, destacando a eficácia desta modificação direcionada~\cite{li2014shift}. O estimador do k-ésimo vizinho mais próximo, integral ao SPEA2, é considerado particularmente propício à integração sinérgica com a metodologia SDE para desempenho aprimorado em muitos objetivos~\cite{li2014shift}.

    Avanços no método de Fertilização In Vitro e aumentações algorítmicas correlatas, como MFO e SDE, mostram promessa para alcançar resultados superiores em otimização multiobjetivo e de muitos objetivos. A confluência do projeto arquitetônico modular (como visto em hibridizações baseadas em IVF), estratégias de manipulação genética e adaptabilidade a domínios de problemas torna essas abordagens de hibridização opções atraentes dentro do conjunto de ferramentas de computação evolutiva. Concomitantemente, frameworks como MFO e técnicas exemplificadas por SDE fornecem mecanismos potentes para amplificar as capacidades operacionais de algoritmos bem estabelecidos, incluindo SPEA2, equipando-os assim para abordar efetivamente paisagens de otimização multiobjetivo novas e formidáveis.